

\title{The $(3,4,\infty)$ family: periods, moduli and the two ends}\author{AI-produced mathematical text}\date{19 November 2026}\maketitle
\paragraph{AI production and check provenance.}
This paper is AI-produced. No human has checked its mathematics.
The recorded stages are GPT-5.6 Sol (OpenAI), in Codex at reasoning
effort ultra, 24 August--9 September 2026, with some contributions by
GPT-5.6 Sol (OpenAI) in ChatGPT Pro mode; GPT-6 Astra (OpenAI), in
ChatGPT Pro mode at effort Pro, for the research runs of 12--16
November, S6\_09 of 16--17 November, the continuing S6\_06 re-derivation
of 16--18 November, and this writing of 17--19 November 2026;
and Claude Fable 5.2 (Anthropic)
at effort high or max and Claude Opus 5.5 (Anthropic) at effort max
for the readings, checks and plans of 12--18 November 2026.
S6\_06\_BLIND is credited as a re-derivation with access to preceding
material. Each statement records the actual check scope and preserves
the names, effort settings and dates of its historical checks.

\paragraph{Effort names.}
The archive's literal settings are retained with their model, provider,
interface and date. OpenAI documents reasoning allocation controls
\cite[sections ``Reasoning effort'' and ``Reasoning mode'']{OpenAIReasoning2026}
and maps ChatGPT's Pro selection to GPT-6 Pro, powered by GPT-6 Astra
\cite[section ``Choose a thinking level'']{OpenAIChatGPTPro2026}.
Anthropic documents effort as control of response-token use, including
thinking \cite[sections ``How effort works'' and ``Effort levels'']{AnthropicEffort2026}.
The check event states what was read, re-derived or executed.

\paragraph{Public-domain dedication.}
This AI-produced text is dedicated to the public domain under
CC0 1.0 Universal \cite{CC01Universal}.


\section{Results and their mechanisms}
\label{sec:I:opening}

\RungZero{We start from the compact complex six-spheres constructed by Claude
and L.~Alp\"oge and Engel's elliptic-surface description: in Paper~I,
theta transformations and integral root norms determine the periods
and moduli, while corner meridians and the Albanese nome obstruct
extension at the two ends; in Paper~II, the torsion term in the
differential sequence identifies the missing vanishing step and gives
the exact replacement criterion; in Paper~III, actual fibre framings
and integral attaching maps determine the homotopy classes and
resulting manifolds. In Paper~IV, the full cusp charts, tensor-frame
return and relative Pontryagin calculation determine the smooth
monodromy on every retained fan, with an explicit family giving a
generator of $\Theta_7$; in Paper~V, rank-six lattice data and their
period extension give the complex fourfold on
$3\mathbb H\mathbb P^2\#(-\mathbb H\mathbb P^2)$; in Paper~VI,
normalization, local Milnor fibres and compact-stratum cohomology
give the Euler sign for all degenerations in its stated
torus-fibration class.}

The fixed data are the period solution with
$(\beta^{\rm part}+\tau)(p_0)=0$, the ordered lattice
$\Lambda=\Z\langle\widehat\gamma,\widehat u,\widehat w,\widehat\delta\rangle$,
clockwise finite root generators, the specified periodic cusp fan, and
the stated affine translations. We write $d=\widehat\delta$,
$E=d\otimes\gamma$, $w=\exp(-\pi i c)$ and $u=w^2$.
The map $f:X\to\Pp^1_t$ has special values $0,1,\infty$.
The construction and its smooth six-sphere are those of Claude and
L.~Alp\"oge, \cite[Theorems~6.2 and 8.1, pp.~36 and 63--64]{AC2026},
and Engel's elliptic-surface description. The contradiction with
\cite[Corollary~2.3, p.~680]{CDP2020} is credited to Claude and
L.~Alp\"oge, \cite[Section~1.2, p.~5; Theorem~10.5 and
Section~10.5, Lemma~10.7, pp.~90--91]{AC2026}; Paper~II gives the exact
replacement for the vanishing step.

\subsection{The period and its exact domain}
\RungZero{For the zero $[\mu]$ of $\wp_\tau$ selected by the cusp
branch, the missing period is
\[
 \beta^{\rm part}=\frac{6i}{\pi}\Log
  \frac{\vartheta_1(\mu,\tau)}{2\sqrt3\,\eta_D(\tau)}.
\]
The two theta transformation identities give the prescribed affine
period laws, and the cusp value fixes the additive constant
(\ref{fam:src:C.1.1}).}
\RungOne
\begin{enumerate}[label=\arabic*.]
\item The pointed elliptic-curve map sends $[\mu]$ to the signed
section with $x=0$. The section has height $1/6$, so the
Mordell--Weil comparison identifies the required generator
(\ref{fam:src:B.1.1}, \ref{fam:src:B.1.3}).
\item The quotient $\vartheta_1(\mu,\tau)/\eta_D(\tau)$ never vanishes:
$\wp$ has its pole at the lattice and its zeros away from that lattice.
It therefore has the specified logarithm on the simply connected
period cover.
\item The transformed quotients acquire the factors
$e^{-\pi i/3+\pi i(1-\mu)^2/\tau}$ and
$i e^{\pi i\mu^2/\tau}$. Multiplication of their logarithms by
$6i/\pi$ gives the two affine laws, with the integer defects fixed at
the finite fixed points.
\item The Fourier equation at the cusp gives
$e^{2\pi i\mu_0}=-(5+2\sqrt6)$ and
$2\sin(\pi\mu_0)=2\sqrt3$. The product then gives
$\beta^{\rm part}+\tau\to0$.
\item The difference from the prescribed solution descends to a
holomorphic function on $\Pp^1$ and has zero cusp value, hence is zero.
\end{enumerate}

\RungZero{The exact domain is $0<|w|<216$. The normalized theta norm
at a zero of $\wp$ has maximum $\sqrt2$, attained precisely over
$t=1$; the real period determinant therefore stays nonzero precisely
on that disc (\ref{fam:src:C.2.3}).}
\RungOne
\begin{enumerate}[label=\arabic*.]
\item The full real determinant has the nonzero factor
$D_c=\Im\beta_c-6(\Im\mu)^2/\Im\tau$, and
$D_{\rm part}=(6/\pi)\log(N(P)/(2\sqrt3))$
(\ref{fam:src:B.1.2}).
\item The isogeny distribution identity yields the division formula
\[
N(P)^{n^2}=\frac{N(nP)|\Delta_E|^{(n^2-1)/12}}{|\psi_n(P)|}.
\]
(\ref{ari:C.1.1}, \ref{ari:C.4.2}).
\item The $n=2,4,8$ formulas and their exact polynomial denominators
bound $N(P)$ strictly below $\sqrt2$ on the three annuli
$|t-1|\ge9/10$, $1/28\le|t-1|\le9/10$, and
$1/500\le|t-1|\le1/28$
(\ref{ari:C.8.1}, \ref{ari:C.8.4}).
\item The local theta expansion gives
$\log N-\log\sqrt2\le-\pi|h|^2/2+150|h|^4<0$
away from $h=0$ on the remaining disc. The two applications of
Rouch\'e's theorem cover that entire $t$-disc
(\ref{ari:C.8.2}, \ref{ari:C.8.3}).
\item At $t=1$, $N=\sqrt2$ and
$M=\max D_{\rm part}=-(3/\pi)\log6$.
At the cusp $D_{\rm part}\to-\infty$. Thus $D_c$ has no zero exactly
when $\Im c<-M$, giving $|w|<e^{-\pi M}=216$.
\end{enumerate}

\subsection{All discrete data and the parameter}
\RungZero{For
$w_1=r_1\varepsilon+k_1d$, $w_2=r_2\varepsilon'+k_2d$,
$3\nmid r_1$, $2\nmid r_2$, and cusp winding $l_0$, the complete
invariants are
\[
 \left(|12l_0-4r_1-3r_2|,\ e^{-\pi i(c+k_2)}\right).
\]
The finite norm equations determine the parity obstruction and the
remaining cusp winding determines the integer $p$
(\ref{fam:src:B.2.1}).}
\RungOne
\begin{enumerate}[label=\arabic*.]
\item The anticanonical pencil and the three labelled special fibres
recover the base coordinate. A biholomorphism consequently intertwines
the two monodromy matrices over that fixed base.
\item Their rational centralizer is $\Q I+\Q E$. The period equation
$L\Pi_c=\Pi_{c'}H$ forces $L=aI_2$, $a=\pm1$, and
$c'-c=n\in\Z$ (\ref{fam:src:A.2.1}).
\item On the root covers the exact equation is
$N_j\ell_j=-(I+nE)w'_j+aw_j$.
The norm images are $3\Z\varepsilon+\Z d$ and
$4\Z\varepsilon'+2\Z d$; hence the last finite obstruction is
$n\equiv k_2-k'_2\pmod2$.
\item A cusp extension has the actual character map
$(x,q)\mapsto(q^v a(q)x^a,q)$, with $v\in\Z^2$ and a holomorphic
unit $a(q)$. Its winding agrees with the two finite defects exactly
when $p'=ap$.
\item The differences of these local maps form a torsor under
$0\to\OO\to\mathcal V\to\OO(-1)\to0$.
Its $H^1$ vanishes, so the displayed local maps and their inverses
glue. This proves sufficiency while retaining the complete collar
section (\ref{fam:src:A.3.1}, \ref{fam:src:B.2.1}).
\end{enumerate}

\RungZero{At each fixed nonzero $u=w^2$ there are two sphere
isomorphism classes, exchanged by a positive $u$ circuit. The
order-four norm quotient detects their difference, and the family
with fixed logarithmic data closes over the punctured $w$-disc
(\ref{fam:src:B.2.2}, \ref{fam:src:B.2.3}).}
\RungOne
\begin{enumerate}[label=\arabic*.]
\item The datum $|p|=1$ is carried to the canonical one by the
explicit finite and cusp maps of \ref{fam:src:B.2.1}.
\item Its remaining identification subgroup is $2\Z$, because
$d$ has order two in $\Lambda^{A_2}/N_2\Lambda$.
\item The logarithms at fixed $u$ differ by integers, producing
exactly the classes $c+2\Z$ and $c+1+2\Z$.
\item A positive circuit sends $c$ to $c-1$, and exchanges these
classes. A descent datum would have to identify the corresponding
fibres; the computed order-two obstruction excludes that datum.
\item The even-shift maps satisfy their exact cocycle law and give
the proper family over $w$. Local inverse logarithms give the
corresponding families on simply connected $u$-open sets.
\end{enumerate}

\RungZero{For coprime multiplicities $(3k,4l)$, the fixed-data period
subgroup is $2kl\Z$, giving $2kl$ classes at fixed $u$.
The two norm images are multiplied by $k$ and $l$, so their central
congruences become $k\mid n$ and $2l\mid n$
(\ref{fam:src:B.3.1}, \ref{push:src:E.1.2}).}
\RungOne
\begin{enumerate}[label=\arabic*.]
\item The vanishing of $H^2(S^6;\R)$ forces every map to
$\Pp^1$ to be constant on a regular K\"ahler torus fibre by the
integral of its pulled-back Fubini--Study form. A biholomorphism
therefore preserves the base and its multiplicities.
\item The root actions have the same linear parts and orders
$3k,4l$, giving $N_{3k}=kN_1$ and $N_{4l}=lN_2$.
\item In the primitive sphere data, the finite equations impose
$k\mid n$ and $2l\mid n$. Their least common multiple is $2kl$
because $\gcd(k,2l)=1$.
\item The cusp equation is $p'=ap$ with
$p=12kl\,l_0-4l\,r_1-3k\,r_2$. Its integral adjustment and the
same vanishing of $H^1(\mathcal V)$ give every allowed comparison.
\item The Chinese remainder map to the two central congruence
groups has $2kl$ elements, and $c\mapsto c-1$ acts by a generator.
\end{enumerate}

\subsection{Division values and the Weierstrass comparison}
\RungZero{At every torsion specialization of $P$, the normalized
norm is algebraic with an exact division-polynomial value,
including the five values below and the complete orders four
through eight: the distribution identity fixes the multiplicative
constant, and cancellation at $mP=\pm P$ gives the algebraic
expression (\ref{ari:C.4.2}, \ref{ari:C.4.3}, \ref{ari:C.4.4}).
\[
\begin{array}{c|c|c}
t&\operatorname{ord}(P)&N(P)^{12}\\ \hline
0&3&27\\
1&2&64\\
32/5&4&10\\
64(47-27\sqrt3)/11&6&17+10\sqrt3\\
64(47+27\sqrt3)/11&6&10\sqrt3-17
\end{array}
\]
The orders at $t=0,1$ are those on the smooth models after
the specified base changes.}
\RungOne
\begin{enumerate}[label=\arabic*.]
\item The theta Green function has zero Haar mean and expansion
$U_\Lambda(z)=\log|z|+\log|\Delta_\Lambda|/12+o(1)$.
These two normalizations determine the isogeny distribution identity.
\item The divisor and leading coefficient of $\psi_m$ cancel the
singularities of $m^2U(z)-U(mz)+\log|\psi_m(z)|$.
The remaining harmonic function on the compact torus is constant,
and the expansion computes that constant.
\item The exact model map
$(x,y)\mapsto(2x,-2\sqrt2y)$ carries the arithmetic cubic and its
section to the geometric cubic and its selected section.
Its discriminant scales by $64$, while the division polynomial
scales by $(-\sqrt2)^{m^2-1}$; the displayed norm identity is
therefore preserved (\ref{I:weierstrass-comparison}).
\item At exact order $n$ choose $m\equiv\pm1\pmod n$.
Then $N(mP)=N(P)$, and
$N(P)^{12}=|\Delta_E|/|\psi_m(P)|^{12/(m^2-1)}$.
\item The nonzero polynomial $\psi_n(P_t)$ makes $t$ algebraic.
The positive radical in the preceding formula proves algebraicity
of $N(P)$ and the stated field inclusions.
\item Evaluate the exact division recurrence at the selected section
and remove the proper-order factors. The resulting squarefree
polynomials on the smooth base are
\[
\begin{aligned}
F_4&=5t-32,\\
F_5&=89t^2+4736t-4096,\\
F_6&=11t^2-6016t+8192,\\
F_7&=881t^4+4745344t^3-12079104t^2\\
   &\qquad+24641536t-16777216,\\
F_8&=49t^4-2614336t^3-7864320t^2\\
   &\qquad+19398656t-8388608.
\end{aligned}
\]
Their exact gcd computations with their derivatives and the
proper-order factors show that their roots give every point of
exact order $4,5,6,7,8$, including the nonreal order-seven roots
(\ref{ari:C.4.4}).
\item For a root of $F_n$, put $r=N(P)^{12}$. Substitution of
$m=n-1$ into the norm identity gives the positive value determined by
\[
\begin{array}{c|l}
n&\text{equation for }r\\ \hline
4&r=64/|t|\\
5&r^5=432^5/((1-t)(5t-32)^4)\\
6&r^2=1728^2/F_5(t)\\
7&r^{35}=2^{192}3^{81}/(|t|^8|t-1|^3|F_6(t)|^{12})\\
8&r^4=1728^4/|F_7(t)|.
\end{array}
\]
At order six the exact identity
$F_5(64(47\mp27\sqrt3)/11)(17\pm10\sqrt3)^2=1728^2$
selects the two positive radicals in the opening table;
$64/(32/5)=10$ gives order four, and the two central theta
values give the remaining entries
(\ref{ari:C.1.2}, \ref{ari:C.4.4}).
\end{enumerate}

\subsection{The multiple fibres and real structures}
\RungZero{The Albanese nomes of the two reduced multiple fibres are
$q_2=iw/216$ and $q_1=-w^2/48^3=(27/64)q_2^2$.
The invariant-coordinate lattice, including its affine half-period,
and the two exact theta values determine these constants
(\ref{ari:C.2.1}).}
\RungOne
\begin{enumerate}[label=\arabic*.]
\item On the order-four product cover the fixed-coordinate lattice
is $\frac12\Z+\frac{\kappa_2}{4}\Z$. Multiplication by two gives
the Albanese lattice $\Z+\frac{\kappa_2}{2}\Z$.
\item The corresponding order-three calculation gives the lattice
with positive-imaginary modulus $-\kappa_1$; both calculations retain
the finite affine deck translations (\ref{fam:src:A.4.1}).
\item Evaluation of the theta expression gives
$\kappa_1=c+\frac12-\frac{3i}{2\pi}\log48$ and
$\kappa_2=c+\frac32-\frac{3i}{\pi}\log6$.
\item Exponentiation in the two displayed Albanese moduli gives
the claimed nomes and their exact relation. Thus $q_2$ identifies
the parameter domain with the punctured unit disc.
\end{enumerate}

\RungZero{The conjugate member has parameter $1-\bar c$, so real
structures occur precisely on the imaginary $w$-diameter, in two
holomorphic-conjugacy classes. The explicit reflection and cap
meridians show that every nonempty fixed manifold is $S^3$
(\ref{fam:src:B.4.2}, \ref{mapsI:src:A.3.1},
\ref{mapsI:src:A.3.2}, \ref{mapsI:src:A.3.3}).}
\RungOne
\begin{enumerate}[label=\arabic*.]
\item The integral matrix
$H=\left(\begin{smallmatrix}1&0&0&0\\0&1&0&0\\-6&0&-1&0\\0&-1&0&-1\end{smallmatrix}\right)$
intertwines the reflected period matrices. The finite affine defects
and the complete cusp unit shift its target from $-\bar c$ to
$1-\bar c$.
\item The moduli invariant sends that parameter to $-\bar w$.
Equality with $w$ is therefore exactly $w\in i\R$ in the punctured
domain. The displayed translations in \ref{mapsI:real-eq-2} make
the reflection square an integral period.
\item Its action on $\operatorname{Aut}(X)=\C^*$ is
$z\mapsto1/\bar z$. The involution equation gives $z\in\R^*$,
and conjugation multiplies $z$ by $|y|^2$; precisely the two signs
remain.
\item The equations $(I-R_\bullet)x=b+\epsilon d/2\pmod\Lambda$
and the cusp equations $z\bar z=K_v$ give every entry of the fixed
fibre table. One negative-imaginary real form maps onto all
$\RP^1$; the other nonempty images are the two specified closed arcs.
\item On each negative-imaginary real form, the semifree circle
action has a fixed circle. Smith homology makes its orbit surface
a disc, and the explicit attachment is
$(D^2\times S^1)\cup(S^1\times D^2)=S^3$.
\item On the positive-imaginary nonempty real form the two cap
meridians, transported by the actual odd-root matrix, have
intersection one. Their genus-one Heegaard gluing gives $S^3$;
the free circle action has Euler number of absolute value one.
\end{enumerate}

\subsection{Integral monodromy and the period domain}
\RungZero{The monodromy group has index twelve in the integral
stabilizer of $d$. It is the kernel of the character
$\chi(B,x,t)=t+6(x_1+x_2+x_1x_2)-\alpha(B)\pmod{12}$;
the commutator and its two conjugates generate the required
Heisenberg lattice (\ref{fam:src:F.1.1}).}
\RungOne
\begin{enumerate}[label=\arabic*.]
\item In the displayed symplectic integral basis every stabilizer
matrix has the unique form $g(B,x,t)$ of
\ref{fam:src:F.1.1}, with $B\in\operatorname{SL}_2(\Z)$,
$x\in\Z^2$, $t\in\Z$.
\item Multiplication gives
$(B,x,t)(C,y,s)=(BC,x+By,t+s-6x^tJ_2By)$.
Substitution proves that $\chi$ is a homomorphism, using the
abelianization character $\alpha(S)=9$, $\alpha(T)=1$.
\item Both geometric generators have $\chi=0$ and their Levi
projections generate $\operatorname{SL}_2(\Z)$.
\item Their commutator word and conjugates supply
$U=(1,0,-6)$, $V=(0,1,-6)$ and $Z=(0,0,12)$ in the unipotent
kernel, with $[U,V]=Z^{-1}$. These generate exactly the
integral solutions of $t+6(x_1+x_2+x_1x_2)=0\pmod{12}$.
\item Reducing any element of $\ker\chi$ first in its Levi
projection and then by $U,V,Z$ puts it in the geometric group.
The central coordinate maps onto $\Z/12$, proving the index.
\end{enumerate}

\RungZero{On the null locus of $\psi(d,Jd)$ there is a dense set of points
with infinite monodromy stabilizer, because the rational parabolic
conjugates of the explicit unipotent stabilizer are dense there.
On either nonzero locus the real parabolic stabilizer is compact,
and the family period map embeds the $t$-line properly in the
quotient: the modular quotient $j/1728$ is its left inverse and
controls the inverse images of compact sets
(\ref{fam:src:F.2.2}, \ref{fam:src:F.2.3}, \ref{fam:src:F.3.1}).}
\RungOne
\begin{enumerate}[label=\arabic*.]
\item The explicit null complex structure is fixed by an infinite
unipotent subgroup. Rational parabolic conjugates remain rational;
the index-twelve lattice calculation supplies an infinite subgroup
of the geometric monodromy in each such stabilizer.
\item Rational parabolic points are dense in the real parabolic
orbit, which is the null locus. These stabilizers give the stated
failure of properness there.
\item At nonzero sign the parabolic is transitive with stabilizer
isomorphic to a compact product. A discrete subgroup therefore acts
properly on that homogeneous space.
\item In the period coordinates the sign invariant is exactly
$\psi(d,Jd)=-6/D_c$. The modular function of the associated
two-dimensional quotient satisfies $j=1728t$.
\item The function $j/1728$ descends to the monodromy quotient
and is a left inverse to the period map. It proves injectivity
and immersion, including the finite root charts. The inverse
image of a compact subset has bounded $t$ and is closed, proving
properness.
\end{enumerate}

\RungZero{The N\'eron--Severi lattice of every regular fibre is
the integral kernel of one displayed period equation, and its
intersection square is $-2\operatorname{Pf}(E)$.
The same Hermitian form gives the projective and elliptic-quotient
criteria with the actual fibre orientation
(\ref{fam:src:D.1.4}, \ref{fam:src:D.1.6}).}
\RungOne
\begin{enumerate}[label=\arabic*.]
\item An integral alternating form $E$ has type $(1,1)$ exactly
when its restriction to the period kernel vanishes. Substitution of
the two kernel vectors gives
$E_{12}-E_{13}\tau+(6E_{23}-E_{14})\mu+E_{24}\beta
+E_{34}(6\mu^2-\tau\beta)=0$.
\item The determinant of the real period map has sign $D_c<0$.
Transporting the integral wedge square through this actual map
therefore gives $E^2=-2\operatorname{Pf}(E)$.
\item The associated Hermitian form is definite exactly when
$\operatorname{Pf}(E)<0$; reversing its sign then gives a positive
integral form and the projective embedding criterion.
\item A nonzero rank-one semipositive integral form has rational
complex kernel. Quotienting by that kernel gives the elliptic
quotient; pulling back an elliptic polarization proves the reverse
implication. Rank one is the equation $\operatorname{Pf}(E)=0$.
\item The monodromy differences of the period functions show that
the generic integral kernel is $\Z\eta$, with $\eta^2=-12$.
The explicit special integral forms in \ref{fam:src:D.1.5}
produce projective fibres for every fixed ordinary $t$.
\end{enumerate}

\subsection{Cohomology and deformations}
\RungZero{The Hodge numbers are the following table, with rows
indexed by $q$ from $3$ to $0$ and columns by $p$ from $0$ to $3$:
\[
\begin{pmatrix}0&0&0&1\\0&1&2&1\\1&2&1&0\\1&0&0&0\end{pmatrix}.
\]
The cotangent direct images reduce them to line bundles on
$\Pp^1$; the last connecting rank is fixed by the nonzero scalar
$\mu_s(z_2)=4i\sqrt\pi/(3^{3/4}\Gamma(1/4)^2)$
(\ref{inv:A2.2}, \ref{inv:A3.1}).}
\RungOne
\begin{enumerate}[label=\arabic*.]
\item The logarithmic relative cotangent sequence is
$0\to N\to E\to\OO\to0$, with
$N=\omega_f(-2S_1-3S_2)$. The invariant logarithmic vertical
field has no zero in that locally free sequence.
\item Relative duality and the root characters compute all
direct images of $N$ and $E$; the extension groups of the
resulting line bundles on $\Pp^1$ vanish
(\ref{inv:A1.1}, \ref{inv:A1.2}).
\item The determinant of the Kodaira--Spencer matrix descends
with simple zero at each finite special value. The theta
half-period derivative evaluates the remaining order-four
entry as the displayed nonzero scalar.
\item The residue sequence, retaining its local torsion
quotients, gives the direct images of $\Omega_X^1$;
duality gives those of $\Omega_X^2$
(\ref{inv:A1.3}, \ref{inv:A2.4}).
\item Leray over the curve has only its two base columns.
The cohomology of the computed line bundles gives precisely
the displayed table and its Serre-dual entries.
\end{enumerate}

\RungZero{The Fr\"olicher sequence satisfies
$E_1\ne E_2=E_\infty$, and the positive-degree ideal in the
Bott--Chern algebra has square zero. The Hodge table and the
Betti numbers force the differential ranks; the fibre class
and its explicit primitives determine the products
(\ref{inv:C1.1}, \ref{inv:C2.1}, \ref{inv:C3.3}).}
\RungOne
\begin{enumerate}[label=\arabic*.]
\item The vanishing of intermediate Betti numbers, together
with the displayed Hodge table, forces the four nonzero
$d_1$ ranks to be one. Every possible later arrow has zero
source or zero target.
\item The resulting double-complex summands determine the
Bott--Chern and Aeppli dimensions through the cited
decomposition theorem, with conjugation and Serre duality
fixing both diagonals.
\item The global Dolbeault representative $Q$ restricts to
the invariant fibre class $\eta$, with $\eta^2=-12$.
Its nonzero restrictions identify the required products
before passage to Bott--Chern cohomology.
\item The displayed $\partial\bar\partial$ primitives in
\ref{inv:C3.3} kill the positive-degree products. The
remaining top-degree candidate has zero integral by Stokes,
which determines its class. Thus the algebra is $\C\oplus N$
with $N^2=0$.
\end{enumerate}

\RungZero{The relative one-forms have torsion exactly along
$S_1\cup S_2$; the cusp module is torsion-free and fails to be
locally free. The explicit modules for $t=s^m$, $t=xy$ and
$t=xyz$ give this correction to the printed torsion description
in Claude--Alp\"oge, Remarks~1.3 and 9.21
(\ref{inv:A4.1}, \ref{inv:A4.2}).}
\RungOne
\begin{enumerate}[label=\arabic*.]
\item At $t=s^m$, the relation is $ms^{m-1}ds=0$,
giving the torsion summand $\OO/(s^{m-1})\,ds$.
\item At $t=xy$, the relation vector $(y,x,0)$ is
primitive in codimension one. Its quotient embeds in the
explicit ideal presentation and has no torsion.
\item At $t=xyz$, the same calculation with
$(yz,xz,xy)$ gives a torsion-free module of rank two
whose fibre dimension jumps on the double locus.
\item Taking the exterior square retains the separate
triple-point torsion module
$(Ru\oplus Rv)/(xu,y(u-v),zv)$. This records exactly
which differential degree each local conclusion concerns.
\end{enumerate}

\RungZero{Every tangent twist has an explicit cohomology
formula. Integral twists are root-line bundles on $\Pp(3,4)$;
nonintegral twists reduce to finite Koszul modules at the
resonant fibres and a length-one cusp term
(\ref{inv:D4.1}, \ref{inv:D4.4}, \ref{inv:D4.5}).}
\RungOne
\begin{enumerate}[label=\arabic*.]
\item For the root generator $\mathcal R$, write
$r_1=n\bmod3$, $r_2=3n\bmod4$ and
$d(n)=(n-4r_1-3r_2)/12$. The coarse invariant sections
of $\mathcal R^n$ form $\OO(d(n))$.
\item The logarithmic tangent Koszul maps have matrices
$(0,y,y,x)^t$ and
$\left(\begin{smallmatrix}-y&0\\-x&y\end{smallmatrix}\right)$.
Their cokernels, followed by the cusp normal quotient,
give the three integral direct images printed in
\ref{inv:D4.1}.
\item A nonintegral twist has nontrivial cusp character;
the normalization complex is acyclic there except for
the stated normal length-one contribution.
\item At a regular resonance the finite relative Dolbeault
complex has $H^1=(A/s)^2$, $H^2=A/s^2$. The root-three
and root-four matrices in \ref{inv:D4.4}, followed by
their exact character invariants, give the other lengths.
\item Coherence makes the resonance support finite for each
twist. Summing its lengths gives $J_z$ and hence
$h^q(T_X\otimes M_z)=(0,J_z,1+J_z,0)$ for nonintegral $z$;
Leray gives the integral formula from the line bundles.
\end{enumerate}

\RungZero{For the standard data, $h^q(T_X)=(1,1,1,0)$, the Kuranishi germ
is smooth of dimension one, and $\operatorname{Aut}(X)=\C^*$:
these credited results of DGKX version~1 are re-derived here through
the root-curve Koszul complex and the integral translation sheaf.
The same root calculation gives $(1,1,1,0)$ for every coprime
$(3k,4l)$ sphere family (\ref{fam:st:A.5.2},
\ref{fam:src:B.4.1}, \ref{push:src:E.2.1}).}
\RungOne
\begin{enumerate}[label=\arabic*.]
\item The weighted root line satisfies
$K_{\Pp(3,4)}=\mathcal R^{-7}$, and its cohomology is
the count of monomials of weights $3$ and $4$.
\item The explicit Kodaira--Spencer matrix and the cusp
normal exact sequence yield the tangent direct images;
their coarse cohomology is $(1,1,1,0)$.
\item Differentiation of the actual period family in $c$
gives the nonzero class $[\xi\otimes\gamma]$.
The induced map to the one-dimensional Kuranishi germ
has nonzero derivative, forcing its defining ideal to
vanish.
\item Every automorphism preserves the recovered base.
The centralizer and finite affine equations force its
linear part to be the identity. The global translation
space modulo its integral periods is $\C e/\Z e=\C^*$.
\item For $(3k,4l)$ the root characters and the same
local Koszul matrices have the invariant degrees
computed in \ref{push:src:E.2.1}; Leray gives the same
four tangent dimensions.
\end{enumerate}
The credited appendices distinguish the two histories. The
working edition contained our proof of the standard-data
$2\Z$ statement by 8 September 2026. DGKX version~1,
dated 23 September 2026, states it as Theorem~1.4;
version~2 does not contain it. Our tangent-cohomology
argument was re-derived after reading DGKX version~1.

\subsection{The puncture and the outer circle}
\RungZero{The untwisted pinching has explicit semistable
and nodal fillings, while its proposed smoothing to the
six-sphere family is excluded by the virtual Euler number.
The plumbing has virtual Euler number zero and a smooth
six-sphere has Euler number two
(\ref{push:src:A.1.2}, \ref{push:src:A.2.2},
\ref{push:src:A.3.2}, \ref{night:P2.10}, \ref{night:P2.11}).}
\RungOne
\begin{enumerate}[label=\arabic*.]
\item The two line-bundle charts are glued by
$z_0w_1=r\varphi$ and $z_1w_0=r$, with the actual unit
$\varphi$ supplied by the translation of the elliptic
section. The annular identifications give a proper family.
\item The blowup of $(z,r)$ and the two small resolutions
of the resulting node give the two triangular fans and
their semistable fillings; contraction gives the nodal one.
\item For the one-component pinching, normalization and the
conductor formula give
$e(Y)-2e(S)+\int_S c_1(S)c_1(L)-c_1(L)^2=24-24=0$.
\item The relative cotangent complex of an arbitrary flat
local-complete-intersection smoothing fixes this virtual
Chern number. On a smooth fibre it is the topological
Euler number, so it excludes the proposed Euler-two
smoothing and gives the correction to Engel's Conjecture~8.2.
\end{enumerate}

\RungZero{After the cover $w=r^6$ the sphere family fills
over $(\Pp^1\times\Delta_r)\setminus\{(\infty,0)\}$,
with a twelve-cycle over the ordinary central base.
The equations $X_iY_i=r$ compactify the long period,
and the actual finite translations permute the components
freely (\ref{push:src:B.4.19}, \ref{night:P2.3}).}
\RungOne
\begin{enumerate}[label=\arabic*.]
\item Put $X_i=W/r^i$, $Y_i=r^{i+1}/W$.
The relations $X_iY_i=r$ and the annular transition
maps give the full infinite chain, including both axes.
\item The long-period deck transformation changes the
index by twelve and translates the elliptic coordinate
by the specified $6\mu$. Its quotient has twelve
central components and compact fibres.
\item On the two finite root charts the component shifts
are $4l_1$ and $3l_2$. Admissibility makes the relevant
cyclic actions free, so the smooth quotients extend there.
\item The cover $y^{12}=t^{4l_1}(t-1)^{3l_2}$ and the
full toric units identify these charts with the original
family for $r\ne0$. The omitted corner is precisely
where both long-period directions meet.
\end{enumerate}

\RungZero{No flat proper filling exists at $w=0$, after any finite cover,
for any admissible twist or any coprime $(3k,4l)$ family.
The actual cusp deck map has $\gamma$-coefficient $-pn$;
commuting meridians propagate the corresponding ratio along the
connected corner fibre, whose two ends give $0$ and $-p\ne0$
(\ref{push:src:B.4.33}, \ref{night:P2.7}, \ref{night:P2.9}).}
\RungOne
\begin{enumerate}[label=\arabic*.]
\item The intrinsic pluricanonical pencil, its labelled
multiple fibres and the complete cusp marking recover
the base map in a putative filling. After modification
and finite cover the two boundary axes are normal crossings.
\item Modulo $\Z d$, a central meridian has finite orbit
under a finite-index subgroup of the base monodromy.
The explicit two unipotent equations have common kernel
zero, so its radial class is zero
(\ref{night:P2.5}, \ref{night:P2.6}).
\item The actual cusp-origin comparison gives
$\Delta_RF=an\,\gamma-6an(n-1)d$, with
$a=-p=4l_1+3l_2-12l_0$. In particular its
$\gamma$-coefficient is $-pn$, with its sign fixed by
the displayed deck maps (\ref{night:P2.2}).
\item For adjacent divisors with orders
$(\alpha_i,\beta_i)$, their commuting coordinate circles
give
$(m\alpha_i+n\beta_i)\kappa_j
=(m\alpha_j+n\beta_j)\kappa_i$.
Thus the ratio is constant across every meeting pair.
\item Properness and connectedness of the corner fibre
join the horizontal divisor to the cusp divisor by
such a chain. The respective ratios are $0$ and $a$,
forcing $p=0$, contrary to admissibility.
\item For $(3k,4l)$ replace twelve by $12kl$ and use
$p=12kl\,l_0-4l\,l_1-3k\,l_2$.
The primitive root equations still make $p\ne0$;
the same connected-chain calculation applies.
\end{enumerate}

\RungZero{No flat proper extension crosses any point of
$|w|=216$. The intrinsic anticanonical divisor carries
an Albanese curve with invariant $J(iw/216)$, whose
rational-cusp growth makes the whole circle a natural
boundary (\ref{push:src:D.1.3}).}
\RungOne
\begin{enumerate}[label=\arabic*.]
\item The unique anticanonical divisor is $2S_2$.
Its reduced support is intrinsic under every fibrewise
biholomorphism, and its Albanese nome is $iw/216$.
\item At a reduced rational $b/d$, modular inversion
gives $\log|j(b/d+iy)|=2\pi/(d^2y)+o(1)$.
This growth excludes a meromorphic continuation there;
density of rational cusps excludes one across any arc.
\item In a putative flat proper extension, remove the
discrete nonsmooth image and the torsion support of
$R^1\pi_*K^{-1}$. The canonical base-change map then
gives a relative anticanonical section of rank one.
\item Its reduced divisor is a smooth proper surface
family away from a countable set. Its Albanese invariant
is holomorphic there, by the explicit weight-one period
construction (\ref{push:src:D.1.2}).
\item An arc has a point outside that countable set.
At that point the invariant continues $J(iw/216)$,
contradicting the growth calculation.
\end{enumerate}

\RungZero{At a boundary parameter the total root quotient
has a Hausdorff central fibre equal to a circle for rational
$a=\Re c+3/2$ and a point for irrational $a$.
All nearby additive-group limits are determined by an
explicit surviving dual lattice; the quadratic term in
the period matrix produces the additional nonseparation
(\ref{push:src:D.2.3}, \ref{push:src:D.3.1}).}
\RungOne
\begin{enumerate}[label=\arabic*.]
\item The integral basis and complex-linear change in
\ref{push:src:D.2.1} give the four periods
$(1,\frac12)$, $(\tau,\frac12+h)$, $(0,1)$,
$(6h,a+B(h))$, where $B(h)=-3ih^2+O(h^4)$.
\item Inverting the transpose of this real period
matrix gives the exact covectors in
\ref{eq:d-dual-exact}. Every bounded sequence has
eventually fixed integral labels.
\item The labels that survive form a primitive subgroup
$K$ of $L_a$; their last coefficients tend to a
homomorphism $\ell:K\to\R$. Pontryagin--Chabauty
duality gives the annihilator formula in
\ref{push:src:D.3.1}.
\item The explicit paths with
$V_n=nv+v_0$ and $d_n(1+\|V_n\|^2)\to0$
realize every such pair $(K,\ell)$ inside the admissible
side. This specifies the maps between the marked limits.
\item On a generic fixed-boundary ray only
$d_0\Re\xi$ survives when $a$ is rational, and no
nonzero label survives when it is irrational. The
resulting period limits prove exactly the asserted
nonseparation relation. Its quotient is Hausdorff and
has the universal factorization property.
\end{enumerate}

\subsection{Arithmetic values and all level curves}
\RungZero{At the eleven rational nonspecial CM parameters,
the theta norm is an explicit positive algebraic factor
times $e^{-\widehat h_O(P)}$. The local division identity
computes every finite-place factor on an integral minimal
model, and the product formula supplies the archimedean
term (\ref{ari:C.6.1}, \ref{ari:C.6.3}).}
\RungOne
\begin{enumerate}[label=\arabic*.]
\item The exact Weierstrass changes in the CM table
send the displayed rational point to the printed
point on each integral minimal model.
\item At each bad prime, the point or its double
lies in the nonsingular identity component. The
valuation of its coordinate and of $\psi_m(P)$ gives
the local height by \ref{ari:C.6.2}.
\item The eleven finite factors are respectively
$\sqrt6$, $2^{2/3}3^{1/4}$, $2^{3/4}$, $2$, $2$,
$\sqrt3$, $1$, $3^{5/12}$, $1$, $1$, $1$ in the
order of the printed CM table.
\item At infinity the same normalization is
$\lambda_\infty(P)=-\log N(P)$.
The product formula for the division identities
therefore gives $N(P)=A_t e^{-\widehat h_O(P)}$.
\item The Mordell--Weil and local calculations in
\ref{ari:C.6.4}--\ref{ari:C.6.5} specify the
non-torsion specializations to which the height
description applies; they do not decide their norm
transcendence.
\end{enumerate}

\RungZero{The norm has one coarse critical point, its
maximum at $t=1$, and each lower level is a Jordan
curve strictly star-shaped about $1$. The exact
derivative is negative along every ray because the
normalized critical factor satisfies $|1-X|<1$
(\ref{ari:C.8.15}, \ref{push:src:D.4.8}).}
\RungOne
\begin{enumerate}[label=\arabic*.]
\item The heat equation gives
$(t-1)L_t=X(X-2)/24$, with
$X=3W\wp'(\mu)/g_2$. The real radial derivative is
$\Re(X(X-2))/12$ (\ref{ari:C.8.12}).
\item The convergent cusp Fourier series gives
$|1-X|<17/20$ when $\Im\tau\ge7/5$
(\ref{ari:C.8.13}).
\item On the remaining Cayley disc, divide out the
known equianharmonic zero. Exact modular recurrences,
Cauchy tails and the displayed positive Bernstein
coefficients give $|1-X|<1$
(\ref{ari:C.8.14}).
\item The two regions cover a closed modular
fundamental domain. The derivative is therefore
strictly negative along each ray from $1$.
The root expansion at $t=0$ gives
$L_t(0)=7/384$, so that point is regular on the
coarse base.
\item The value tends to $\log\sqrt2$ at $1$ and
to $-\infty$ uniformly at infinity. Each ray meets
each lower level exactly once; the nonzero derivative
gives a continuous, locally regular radial graph.
That graph is the asserted Jordan curve.
\end{enumerate}

\subsection{The exact remaining questions}
\RungZero{Two questions remain open in the retained
scope: a literal coefficient atlas for the global
Bott--Chern representatives, and transcendence of the
normalized theta norm at the specified non-torsion
arithmetic points. The known global classes and the
exact height identities reduce these questions to
the concrete gaps recorded in the final part
(\ref{inv:C3.8}, \ref{ari:C.7.1}--\ref{ari:C.7.12}).}
\RungOne
\begin{enumerate}[label=\arabic*.]
\item The sheaf and double-complex calculations
already construct the global classes and determine
their dimensions and products. The remaining atlas
problem asks for the specified coefficient functions
on every finite and cusp chart with the stated
overlap identities.
\item At torsion points the division identity
already determines an algebraic value of the norm.
At the non-torsion points the proved relation is
$N(P)=A_t e^{-\widehat h_O(P)}$, with $A_t$
explicit and algebraic.
\item The period comparisons and the conductor-361
calculation identify the exact further
algebraic-independence input used by the attempted
transcendence route. That input is recorded as an
open step, together with every established identity;
no transcendence conclusion is asserted from it.
\end{enumerate}


